\documentclass[11pt]{article}
\usepackage[utf8]{inputenc}
\usepackage[spanish]{babel}
\usepackage{amsmath, amssymb, amsthm}
\usepackage{geometry}
\usepackage{enumitem}
\usepackage{tcolorbox}
\usepackage{booktabs}
\geometry{margin=1in}

% Colores para cajas
\definecolor{blue1}{RGB}{0,82,155}
\definecolor{red1}{RGB}{186,24,27}
\definecolor{green1}{RGB}{0,128,0}

\title{\textbf{Tutoria 1 (Pauta) - Algebra Lineal}}
\author{Marcelo Alegría}
\date{}

\begin{document}

\maketitle

\section*{Ejercicio 1: Independencia Lineal}

\begin{tcolorbox}[colback=blue1!10, colframe=blue1, title=Enunciado]
Considere los vectores $\vec{v}_{1}=\begin{pmatrix}1\\ 2\\ 3\end{pmatrix}$, 
$\vec{v}_{2}=\begin{pmatrix}2\\ 4\\ 6\end{pmatrix}$, 
$\vec{v}_{3}=\begin{pmatrix}3\\ 6\\ 9\end{pmatrix}$ y 
$\vec{v}_{4}=\begin{pmatrix}1\\ 1\\ 1\end{pmatrix}$.

\begin{enumerate}[label=\alph*)]
    \item ¿Son los vectores linealmente independientes?
    \item ¿Son los vectores linealmente independientes si eliminamos $\vec{v}_{3}$?
    \item ¿Son los vectores linealmente independientes si eliminamos $\vec{v}_{2}$ y $\vec{v}_{3}$?
\end{enumerate}
\end{tcolorbox}

\begin{tcolorbox}[colback=green1!10, colframe=green1, title=Solución]
\textbf{a)} Analizamos la matriz formada por los vectores:
\[
A = \begin{pmatrix}
1 & 2 & 3 & 1 \\
2 & 4 & 6 & 1 \\
3 & 6 & 9 & 1
\end{pmatrix}
\]

Observamos que:
\begin{itemize}
    \item $\vec{v}_2 = 2\vec{v}_1$
    \item $\vec{v}_3 = 3\vec{v}_1$
\end{itemize}
Por lo tanto, $\{\vec{v}_1, \vec{v}_2, \vec{v}_3\}$ son linealmente dependientes. 
Aunque $\vec{v}_4$ es LI con $\vec{v}_1$, el conjunto completo es \textbf{LD}.

\textbf{b)} Eliminando $\vec{v}_3$: $\{\vec{v}_1, \vec{v}_2, \vec{v}_4\}$
\[
\begin{vmatrix}
1 & 2 & 1 \\
2 & 4 & 1 \\
3 & 6 & 1
\end{vmatrix} = 0 \quad \text{(primeras dos columnas proporcionales)}
\]
Sigue siendo \textbf{LD} porque $\vec{v}_2 = 2\vec{v}_1$.

\textbf{c)} Eliminando $\vec{v}_2$ y $\vec{v}_3$: $\{\vec{v}_1, \vec{v}_4\}$
\[
\text{Rango}\left(\begin{pmatrix}1 & 1\\ 2 & 1\\ 3 & 1\end{pmatrix}\right) = 2
\]
Estos dos vectores son \textbf{LI} (no son múltiplos escalares).
\end{tcolorbox}

\section*{Ejercicio 2: Combinaciones Lineales y Generados}

\begin{tcolorbox}[colback=blue1!10, colframe=blue1, title=Enunciado]
\begin{enumerate}[label=\alph*)]
    \item Escriba el vector $\begin{pmatrix} 3 \\ 4 \\ 5 \end{pmatrix}$ como combinación lineal de 
    $\begin{pmatrix} 1 \\ 0 \\ 1 \end{pmatrix}$, $\begin{pmatrix} 1 \\ 1 \\ 1 \end{pmatrix}$, 
    $\begin{pmatrix} 1 \\ 1 \\ 0 \end{pmatrix}$.
    
    \item Encuentre $\text{gen}\left(\begin{pmatrix} 1 \\ 2 \\ 4 \end{pmatrix},\begin{pmatrix} 1 \\ 0 \\ 0 \end{pmatrix}\right)\cap 
    \text{gen}\left(\begin{pmatrix} 1 \\ 2 \\ 3 \end{pmatrix}\right)$.
\end{enumerate}
\end{tcolorbox}

\begin{tcolorbox}[colback=green1!10, colframe=green1, title=Solución]
\textbf{a)} Resolvemos el sistema:
\[
x\begin{pmatrix}1\\0\\1\end{pmatrix} + y\begin{pmatrix}1\\1\\1\end{pmatrix} + z\begin{pmatrix}1\\1\\0\end{pmatrix} = \begin{pmatrix}3\\4\\5\end{pmatrix}
\]
Sistema de ecuaciones:
\begin{align*}
x + y + z &= 3 \\
y + z &= 4 \\
x + y &= 5
\end{align*}
Resolviendo: $x = 1$, $y = 4$, $z = 0$. Por lo tanto:
\[
\begin{pmatrix}3\\4\\5\end{pmatrix} = 1\cdot\begin{pmatrix}1\\0\\1\end{pmatrix} + 4\cdot\begin{pmatrix}1\\1\\1\end{pmatrix} + 0\cdot\begin{pmatrix}1\\1\\0\end{pmatrix}
\]

\textbf{b)} El primer generado es:
\[
\text{gen}\left\{\begin{pmatrix}1\\2\\4\end{pmatrix},\begin{pmatrix}1\\0\\0\end{pmatrix}\right\} = 
\left\{a\begin{pmatrix}1\\2\\4\end{pmatrix} + b\begin{pmatrix}1\\0\\0\end{pmatrix} \mid a,b\in\mathbb{R}\right\}
= \left\{\begin{pmatrix}a+b\\2a\\4a\end{pmatrix} \mid a,b\in\mathbb{R}\right\}
\]

El segundo generado es:
\[
\text{gen}\left\{\begin{pmatrix}1\\2\\3\end{pmatrix}\right\} = \left\{c\begin{pmatrix}1\\2\\3\end{pmatrix} \mid c\in\mathbb{R}\right\}
= \left\{\begin{pmatrix}c\\2c\\3c\end{pmatrix} \mid c\in\mathbb{R}\right\}
\]

La intersección son los vectores que cumplen:
\[
\begin{pmatrix}a+b\\2a\\4a\end{pmatrix} = \begin{pmatrix}c\\2c\\3c\end{pmatrix}
\Rightarrow
\begin{cases}
a + b = c \\
2a = 2c \\
4a = 3c
\end{cases}
\Rightarrow
\begin{cases}
a = c \\
4a = 3c \\
a + b = c
\end{cases}
\Rightarrow a = c = 0, b = 0
\]

Por lo tanto, la intersección es solo el vector nulo: $\{\vec{0}\}$.
\end{tcolorbox}

\section*{Ejercicio 3: Sistemas con Parámetros}

\begin{tcolorbox}[colback=blue1!10, colframe=blue1, title=Enunciado]
Sea $a \in \mathbb{R}$ y considere el sistema cuya matriz aumentada viene dada por:
\[
\begin{pmatrix}
2 & 0 & -3 & | & 0 \\
0 & 2a & 0 & | & -1 \\
0 & 0 & a^2 - 1 & | & a + 1
\end{pmatrix}
\]

\begin{enumerate}[label=\alph*)]
    \item Determine los valores de $a$ para los cuales la matriz anterior está en su forma escalonada.
    \item Determine los valores de $a$ para que el sistema tenga: única solución, infinitas soluciones, ninguna solución.
\end{enumerate}
\end{tcolorbox}

\begin{tcolorbox}[colback=green1!10, colframe=green1, title=Solución]
\textbf{a)} La matriz está en forma escalonada si:
\begin{itemize}
    \item Los pivotes (2, $2a$, $a^2-1$) son diferentes de cero
    \item O si son cero, deben estar en la última fila
\end{itemize}
Condiciones: $2a \neq 0 \Rightarrow a \neq 0$ y $a^2 - 1 \neq 0 \Rightarrow a \neq \pm 1$.

\textbf{b)} Analizamos por casos:

\begin{itemize}
    \item \textbf{Solución única}: Todos los pivotes $\neq 0$
    \[
    a \neq 0 \quad \text{y} \quad a^2 - 1 \neq 0 \Rightarrow a \neq 0, \pm 1
    \]
    
    \item \textbf{Infinitas soluciones}: Última fila $[0\ 0\ 0\ |\ 0]$
    \[
    a^2 - 1 = 0 \quad \text{y} \quad a + 1 = 0 \Rightarrow a = -1
    \]
    
    \item \textbf{Sin solución}: Última fila $[0\ 0\ 0\ |\ b]$ con $b \neq 0$
    \[
    a^2 - 1 = 0 \quad \text{y} \quad a + 1 \neq 0 \Rightarrow a = 1
    \]
    
    \item \textbf{Caso especial}: $a = 0$
    \[
    \begin{pmatrix}
    2 & 0 & -3 & | & 0 \\
    0 & 0 & 0 & | & -1 \\
    0 & 0 & -1 & | & 1
    \end{pmatrix}
    \Rightarrow \text{Fila 2: } 0 = -1 \Rightarrow \text{No hay solución}
    \]
\end{itemize}

Resumen:
\begin{itemize}
    \item \textbf{Única solución}: $a \neq 0, \pm 1$
    \item \textbf{Infinitas soluciones}: $a = -1$
    \item \textbf{Sin solución}: $a = 0, 1$
\end{itemize}
\end{tcolorbox}

\section*{Ejercicio 4: Aplicación - Sistema de Producción}

\begin{tcolorbox}[colback=blue1!10, colframe=blue1, title=Enunciado]
Una empresa produce tres tipos de productos: $P_1$, $P_2$ y $P_3$. La empresa tiene tres tipos de recursos: material, mano de obra y energía. Las cantidades requeridas son:

\begin{center}
\begin{tabular}{lccc}
\toprule
 & Material & Mano de obra & Energía \\
\midrule
$P_1$ & 3 unidades & 2 horas & 1 unidad \\
$P_2$ & 2 unidades & 4 horas & 2 unidades \\
$P_3$ & 4 unidades & 1 horas & 3 unidades \\
\bottomrule
\end{tabular}
\end{center}

Disponible: 20 unidades de material, 18 horas de mano de obra, 15 unidades de energía.

\begin{enumerate}[label=\alph*)]
    \item Plantee un sistema de ecuaciones para utilizar todos los recursos.
    \item Resuelva el sistema.
\end{enumerate}
\end{tcolorbox}

\begin{tcolorbox}[colback=green1!10, colframe=green1, title=Solución]
\textbf{a)} Sea $x_1$, $x_2$, $x_3$ las unidades de $P_1$, $P_2$, $P_3$ respectivamente.

Sistema de ecuaciones:
\begin{align*}
3x_1 + 2x_2 + 4x_3 &= 20 \quad \text{(material)} \\
2x_1 + 4x_2 + x_3 &= 18 \quad \text{(mano de obra)} \\
x_1 + 2x_2 + 3x_3 &= 15 \quad \text{(energía)}
\end{align*}

\textbf{b)} Resolvemos por eliminación gaussiana:

Matriz aumentada:
\[
\begin{pmatrix}
3 & 2 & 4 & | & 20 \\
2 & 4 & 1 & | & 18 \\
1 & 2 & 3 & | & 15
\end{pmatrix}
\]

Operaciones:
\begin{align*}
& F_1 \leftrightarrow F_3: \begin{pmatrix}1 & 2 & 3 & | & 15 \\ 2 & 4 & 1 & | & 18 \\ 3 & 2 & 4 & | & 20\end{pmatrix} \\
& F_2 \leftarrow F_2 - 2F_1: \begin{pmatrix}1 & 2 & 3 & | & 15 \\ 0 & 0 & -5 & | & -12 \\ 3 & 2 & 4 & | & 20\end{pmatrix} \\
& F_3 \leftarrow F_3 - 3F_1: \begin{pmatrix}1 & 2 & 3 & | & 15 \\ 0 & 0 & -5 & | & -12 \\ 0 & -4 & -5 & | & -25\end{pmatrix} \\
& F_2 \leftrightarrow F_3: \begin{pmatrix}1 & 2 & 3 & | & 15 \\ 0 & -4 & -5 & | & -25 \\ 0 & 0 & -5 & | & -12\end{pmatrix}
\end{align*}

Sistema escalonado:
\begin{align*}
x_1 + 2x_2 + 3x_3 &= 15 \\
-4x_2 - 5x_3 &= -25 \\
-5x_3 &= -12
\end{align*}

Solución:
\begin{align*}
x_3 &= \frac{12}{5} = 2.4 \\
x_2 &= \frac{25 - 5\cdot 2.4}{4} = \frac{13}{4} = 3.25 \\
x_1 &= 15 - 2\cdot 3.25 - 3\cdot 2.4 = 15 - 6.5 - 7.2 = 1.3
\end{align*}

Por lo tanto: $x_1 = 1.3$, $x_2 = 3.25$, $x_3 = 2.4$ unidades.
\end{tcolorbox}

\end{document}